\section{The Cartier normal form on arbitrary saturated rows}
\label{sec:cartier-normal-v59}

Fix the original row datum $X$, its alternating type pattern, a side
$g$, a big type $a$, and an odd integer $i<d$. Put $p=V.\ell$,
$v(p)=1$, $R=\{v\ge0\}$, $\mathfrak m=\{v>0\}$,
$k=R/\mathfrak m$, $h=(i-1)/2$, $K=h+1$, and $q=p^2$. The field
$C$ is algebraically closed of characteristic zero by the row datum.
The residue field has characteristic $p$. Evenness of $d$ makes parity
well defined in $\mathbb Z/d$. On side $g$ each $a+2s$ is big
and each $a+2s+1$ is small. On the opposite side these designations
are reversed. Write $U,U'$ for the supplied $n$-row matrices, of
types $a$ and $a+i$ respectively. Saturation means integral entries
and an integral right inverse. Row independence modulo $\mathfrak m$
means linear independence over $k$. This follows by lifting the
coefficients of a residue relation to $R$. The weighted conclusions
below assume $p\ge5$.

% BEGIN CRT1
\paragraph{CRT1. Integral elimination on the given Cartier curves.}
Write $x_c(t)$ for the $t$-th vector of a curve and $B_c$ for its
tangent row, so $x_c(0)=(B_c,0)$. Write $[u]x(t)=u^{p^t}x(t)$ and
$Vx(t)=x(t-1)$ for $t>0$, with $Vx(0)=0$. The curve module admits both
operations for $u\in R$. The tangent basis $B$ is integral and
unimodular; consequently $B^{-1}$ is integral. A pure integral
tangent of type $b$ has a curve lift of type $b$ by the defining
Cartier lifting property. Applied row by row, this realizes any pure
integral unimodular matrix as a tangent system.

Here are the coefficient bounds needed for changing this system.
The Frobenius shift $Fx(t)=p x(t+1)$ preserves the curve module.
Hence $v(x(t))\ge-t$, since $F^tx$ has integral initial value.
More generally, if the first $s$ vectors of a module element vanish,
then $v(x(s+t))\ge-t$. Indeed $x(s)$ is integral by the initial
coefficient axiom at $s$. Lift its $I$-part using the pure unit
tangent curves and integral Teichm\"uller combinations, and its
$J$-part by the allowed curve concentrated at zero. Subtract the
$s$-fold shift of this lift. The remainder begins at $s+1$;
induction on $t$ bounds its coefficient there and the unshifted
lift contributes valuation at least $-t$. For a good curve
($x(1)_I=0$), $Fx$ minus the constant curve with $J$-part
$p x(1)_J$ begins at $1$. We obtain
$v(x(t))\ge1-t$ for $t\ge2$. All bounds are entrywise.

Reduce $U$ and its right inverse. The reduction of $U$ has rank
$n$, so choose $n$ pivot columns with unit determinant. They have
type $a$. At these positions replace the corresponding rows of the
identity matrix by the rows of $U$, in pivot order. Its determinant
is the pivot minor, by expansion along the untouched rows. The
preceding lifting property gives a system and a set $Q$ of $n$
marked rows, with tangents in $L=\sum_b R U_b$.

We need integral changes within a marked set $Q$ of one type.
Given one column of numbers $z_c$ on $Q$, retain $Q$ if all vanish.
Otherwise, pivot at an entry of smallest valuation and subtract
$(z_c/z_{c_0})$ times the pivot row from every other marked row.
The transformation lies in $\mathrm{GL}(R)$ and is the identity
outside $Q$. Exclude the pivot from further elimination. Repeating
on a family of $m$ columns preserves established zeros on surviving
rows and loses at most $m$ rows. If the column entries $w_{ce}$ are
integral and the desired operations have exponent $p^t$, first take
$p^t$-th roots of every entry in $C$. These roots are integral.
The preceding operations give a single matrix $P$ for which the
linear sums of the roots vanish on survivors. Modulo $pR$, Frobenius
is additive, so
\[
             \sum_x P_{cx}^{p^t} w_{xe}\in pR .
\]
The ideal here is $pR$, which can be strictly smaller than
$\mathfrak m$. We have used characteristic $p$ of $R/pR$,
without a discreteness assumption. Replacing $x_c$ by
$\sum_x[P_{cx}]x_x$ gives exactly this transformation at
degree $t$, and ordinary $P$ at the tangent. Purity, the determinant
condition, and membership of marked tangents in $L$ are preserved.
% END CRT1

% BEGIN CRT2
\paragraph{CRT2. Constructing the adapted chain.}
Suppose the current marked rows have type $b$. Apply the preceding
elimination at $t=1$ to the integral vectors $p x_c(1)_I$, in the
columns of type $b+1$. After at most the number of such columns are
lost, $x_c(1)_I$ is integral on each survivor and zero elsewhere
in type by purity. Lift the integral tangent $x_c(1)_I$ by a curve $z$ with
$z(0)=(x_c(1)_I,0)$
of type $b+1$ and subtract $Vz$. Now the survivor is good and its
tangent is unchanged.

For a set of good rows of type $b$, write $p x_c(2)_I$ in the
current tangent basis. Its coordinates are integral by CRT1.
Fix a set $E$ of forbidden indices of type $b+2$. Eliminate at
$t=2$ on these columns. This preserves goodness and loses at most
$|E|$ additional rows. Express the resulting scaled vectors using
the tangent basis from immediately before this change. The entries
of their coefficient matrix at $E$ belong to $pR$. The basis is
block diagonal by type, hence the coefficients at other types
vanish. Separate the $E$ terms to obtain on surviving rows
\[
    p x_c(2)_I=\sum_{e\text{ of type }b+2,\ e\notin E}
                 g_{ce}B_e+p\lambda_c,\qquad
                 g_{ce},\lambda_c\in R
\]
(entrywise for $\lambda_c$). In the uses with a next level, these
$B_e$ lie outside the changing type and are its final tangents.

The order of this construction controls the total loss. For a chain
of $K'$ consecutive levels starting at type $b$, first construct the
tail of length $K'-1$ on \emph{all} rows of type $b+2$, using the
entire tangent module as its allowed span. For length zero there are
no conditions. The tail construction misses at most $(K'-1)j'$
of its candidate initial rows. It leaves the current type $b$
untouched since the relevant even offsets are positive and less
than $d$. Now impose goodness at $b$, losing at most $j'$ (the
intermediate type is small). Take $E$ to be the missed rows of
the tail and apply the second elimination. Thus at most $K'j'$
current candidates are lost. For $K'=1$ only the goodness step
is needed. Number the survivors at $b$ by level zero and the
previous levels by their old level plus one; assign unused indices
a label at least $K'$. Changes within the current type preserve every
previously constructed curve in the tail and all of its supported
tangent coefficients. This proves the asserted induction with the
original marked lattice $L$ on the outermost step.

Let $N_1$ be the number of surviving marked rows at zero and let
$\iota$ enumerate them. Then $n\le N_1+Kj'$, in particular
$n\le N_1+(K+1)j'$ as stipulated. Define $D.\alpha_{ctb}$ to be
the $b$-th tangent-basis coefficient of $x_c(t)_I$. At every active
level these curves are good, their type is $a+2\operatorname{lev}(c)$,
and CRT1 supplies all the $-t$ and $1-t$ bounds. The first
coefficient is zero on good rows. For a row whose successor is
active, the preceding display in tangent coordinates gives
$pD.\alpha_{c2e}=g_{ce}+p\lambda_{ce}$ with $g$ supported at the
next level. At the terminal level take $g=pD.\alpha_{c2*}$,
$\lambda=0$, and off the active levels take both zero. This gives
all fields of the original adapted data with valuation $v$, prime
$p$, and index $i$. The zero-level tangents belong to $L$ by the
construction, providing actual integral coefficients $A$ in
$B_{\iota c}=\sum_b A_{cb} U_b$. The selected tangents are
saturated by the corresponding columns of $B^{-1}$.
% END CRT2

% BEGIN CRT3
\paragraph{CRT3. Selection by rank at every level.}
For each level $s$ write $F_s=\{c:\operatorname{lev}(c)=s\}$ and put
\[
 P_{0,cj}=\mathbf1_{c=\iota j},\qquad
 P_{s+1,ej}=\sum_{c\in F_s}pD.\alpha_{c2e}\,P_{s,cj}^{q}.
\]
These are the original chain matrix entries on the initial labels.
They are integral as long as a level is used. At a nonterminal
step the entries off the successor level belong to $pR$.
Suppose integral rows $Y_\delta$ restricted to $F_s$ are independent
over $k$. Their images on $F_{s+1}$ have entries
$Z_{\delta e}=\sum_{c\in F_s}pD.\alpha_{c2e}Y_{\delta c}^{q}$.
If independent residue rows $x_\nu$ satisfy
$\sum_\delta x_{\nu\delta}^{q}Z_{\delta e}=0\pmod{\mathfrak m}$
on that next level, put $w_{\nu c}=\sum_\delta x_{\nu\delta}Y_{\delta c}$.
The Teichm\"uller combination $\sum_{c\in F_s}[w_{\nu c}]x_c$
is good of the common big type $a+2s$. Its scaled second tangent
is zero modulo $\mathfrak m$: expressed in $B$, its coefficients
are $\sum_{c\in F_s} w_{\nu c}^{q}pD.\alpha_{c2e}$, which vanish on
the next level by residue Frobenius and off it by $pR$ divisibility.
Its tangents are independent modulo $\mathfrak m$, by row independence
of $x$, then $Y|_{F_s}$, and invertibility of $B$. The original
row axiom $R1$ gives $\#\{\nu\}+\mathrm{rh}\le j'$. In particular
the loss from this step is at most $j'$.

For clarity, the selection principle even works with any integer
$Q\ge1$ in the powered relation. Choose a maximal independent
subset of the residue rows of an integral matrix. For each excluded
row lift a dependence on this subset with its excluded entry equal to
one and other excluded entries zero. Take entrywise $Q$-th roots in
$C$, choosing $0$ and $1$ on the excluded columns. The rows of roots
are integral and independent there by their identity minor, and their
$Q$-th powers are relations modulo $\mathfrak m$. A bound on the
size of any such independent family bounds the number of excluded
rows. This does not require $z\mapsto z^Q$ to be additive for
arbitrary $Q$. Apply this principle with $Q=q$ at each step,
starting with the distinct standard columns at zero. At level $h$
we have retained an actual subset of $m$ initial indices with
$N_1\le m+h j'$ whose $P_{h,*,j}|_{F_h}$ are independent.
% END CRT3

% BEGIN CRT4
\paragraph{CRT4. Terminal rank, the pairing, and the integral row operation.}
Set $\beta_{ct\kappa}=x_c(t)_{J,\kappa}$ and, for the retained labels, put
\[
 T_{j\kappa}=(-1)^{h+1}\sum_{c\in F_h}P_{h,cj}^{p}\,p\beta_{c1\kappa}.
\]
Each vector is integral and has type $a+i$ by purity. If integral
coefficients $d_j$ give a zero residue combination of the $T_j$,
choose integral $p$-th roots $r_j^p=d_j$ and form the good curve
with Teichm\"uller coefficients $w_c=\sum_j r_jP_{h,cj}$ on $F_h$.
Its scaled first $J$-vector is zero in residue by Frobenius.
The original row axiom $R2$ makes its tangent zero there. Since
$B$ is invertible, $w|_{F_h}=0$ in residue; independence at the
final level forces all $r_j$, hence $d_j$, into $\mathfrak m$.

Write $H$ for the supplied integral pairing matrix from $J_g$ to
$I_{1-g}$ and $E'$ for the coordinate space of opposite type $a+i$.
It has dimension $n+j$. The images $T_jH$ are independent over $k$
because $H$ has an integral right inverse, and they are supported on
$E'$ by type compatibility. The restrictions of $U'$ to $E'$ are
$n$ independent rows by saturation. Hence the residue annihilator
of $U'$ inside $(E')^*$ has dimension $j$. The map from the span
of the $T_j$ given by $T\mapsto TH(U')^{\mathsf t}$ has kernel
of dimension at most $j$. Select $n_0$ of these original labels
with independent images and $m\le n_0+j$. Let $e$ consist of their
tangents and $T$ denote their terminal matrix. Thus $n_0\le n$ and
\[
 n\le n_0+(K+1+h)j'+j=n_0+(i+1)j'+j .
\]
Write $D'=U'H^{\mathsf t}T^{\mathsf t}$. Its $n_0$ columns are
residue independent. Choose standard coordinate columns to complete
them to an $n\times n$ matrix $M$ invertible over $R$, by completing
its residue columns to a basis using coordinate vectors. Set
$Y=-M^{-1}$, $v'=YU'$, $\Phi=YD'$. Then $v'$ is obtained by an
integral unimodular operation, retains type $a+i$, and
$\Phi_{bc}=-\delta_{bc}$ for $b<n_0$. In fact the lower entries
of this construction are also zero. Crucially
$\Phi_{bc}=\sum_{\kappa,k}T_{c\kappa}H_{\kappa k}v'_{bk}$,
using the very same terminal vectors and the original pairing.
The construction includes $n_0=0$ with empty columns.
% END CRT4

% BEGIN CRT5
\paragraph{CRT5. Weighted coordinates on the selected leaf.}
Use the system just constructed on $g$ and choose one with identity
tangent on the other side. Each has unique zero-constant coordinates
$t$: the equation for degree one has invertible tangent matrix and
every higher-degree term involving $t_c^{p^r}$, $r>0$, depends
only on lower degrees. The original logarithm equation therefore
determines the series coefficient by coefficient. Set $z_j$ to be
the left variables and restrict coordinates along the chosen tangent
rows $e_j=B_{u_j}$, where $u_j$ denotes their original index.
Write $L_c$ for the pullback of $t_c$ on this side. In the tangent
basis the defining equations become
\[
 L_b=\sum_{j:u_j=b}z_j-\sum_{r\ge1,c}\alpha_{crb}L_c^{p^r},
 \qquad Z_\kappa=-\sum_{r\ge1,c}\beta_{cr\kappa}L_c^{p^r}.
 \tag{C1}
\]
Here $Z_\kappa$ is the pullback of the original $S.\mathrm{ze}(t)_\kappa$.
The first equation follows by multiplying the coordinate vector by
$B^{-1}$; the second is its defining $J$-series. Each source and
target substitution has zero constant term, so coefficients at a fixed
degree involve just finitely many terms.

Set $\rho=1/(p^2-1)$, $\sigma=3/(8p^i)$, $r_L=\rho+\sigma$.
Write $w(F)\ge b$ when $v([z^\mu]F)+r_L|\mu|\ge b$ for every
$\mu$. Sums preserve lower bounds and products add them, by finite
coefficient convolution. For the \emph{selected} initial labels
now let $P_{s,cj}$ denote the same chain entries with initial column
$u_j$. Define
\[
 D_s=1+q+\cdots+q^{s-1},\quad
 k_s=(-1)^sp^{-D_s},\quad
 A_{s,c}=k_s\sum_j P_{s,cj}z_j^{q^s},\qquad D_0=0.
\]
Each monomial has weight at least $\rho+q^s\sigma$, since
$D_s=(q^s-1)\rho$. For $0\le s\le h$ put $A_c=A_{s,c}$ if
$c\in F_s$, and put $A_c=0$ on the other indices. We prove
\[
                         w(L_c-A_c)\ge1+\rho.             \tag{C2}
\]
Consider the right-hand side of (C1) as a substitution operator on
zero-constant vectors $L'$, and assume $w(L'_c-A_c)\ge1+\rho$.
For $c$ on an active level the $r=1$ summand vanishes. For $r\ge3$
its weight is at least $1-r+p^r\rho\ge1+\rho$; equivalently
$p^r-1\ge r(p^2-1)$ for $r\ge3$, true at $r=3$ for $p\ge3$ and
preserved on incrementing $r$. For an inactive $c$ every $r\ge1$
has weight at least $-r+p^r(1+\rho)\ge1+\rho$.

The exceptional $r=2$ term propagates exactly the pure powers.
For any $p$-power $m$, write $L'_c$ as a sum of the pure monomials
of $A_c$ and the single remainder $E_c=L'_c-A_c$. In its $m$-th
power every mixed multinomial coefficient is divisible by $p$,
as seen by reducing the polynomial identity modulo $p$. Hence the
difference from the pure $m$-th powers of the displayed monomials
of $A_c$ has weight at least
$\min\{1+m\rho,m(1+\rho)\}$. For $m=q$ and the factor
$\alpha_{c2b}$ of valuation at least $-1$, this leaves at least
$1+\rho$ since $q\rho=1+\rho$. The remaining pure contributions,
including the minus sign in (C1), sum on $F_s$ to $A_{s+1,b}$:
this is the recurrence for $P$ and $k_{s+1}=-k_s^q/p$.
At each nonterminal step, coefficients landing off their intended
level gain one in valuation by CRT2. The final escape $A_{h+1,b}$
has weight at least $\rho+q^{h+1}\sigma=\rho+3p/8\ge1+\rho$.
The linear term is $A_{0,b}$, zero off $F_0$. These estimates prove
that the operator carries the set $w(L'-A)\ge1+\rho$ into itself.
Start at $A$ and iterate. Agreement of series through degree $m$
implies agreement of their images through degree $m+1$, since the
nonlinear powers have exponent at least two and all constants vanish.
The stabilized coefficients solve (C1), proving (C2) even without
completeness of $C$. In particular $w(L_c)\ge0$ for every $c$.

For $\kappa$ of type $a+i$, the term $\beta_{c1\kappa}$ at every
earlier active level is zero: it has type $a+2s+1$, and
$0< i-(2s+1)<d$ there. For active $c$ and $r\ge2$ the summand
in $Z_\kappa$ has weight at least $1-r+p^r\rho\ge0$, since
$p^{r-2}\ge r-1$. On inactive rows the bound is
$-r+p^r(1+\rho)\ge0$. At $r=1$ on $F_h$, the preceding mixed-term
estimate with $m=p$ leaves nonnegative weight after multiplication
by $\beta_{c1\kappa}$ of valuation at least $-1$. Its pure terms
give precisely
\[
 w\!\left(Z_\kappa-p^{-E}\sum_j T_{j\kappa}z_j^{p^i}\right)\ge0,
 \qquad E=1+pD_h .                                            \tag{C3}
\]
Indeed the sign is $-(-1)^{hp}=(-1)^{h+1}$, the exponent is
$pq^h=p^i$, and $p\beta_{c1\kappa}$ contributes the extra
denominator $p$. Equations (C2)--(C3) use the same selected indices,
curves and terminal matrix as CRT4. Along with their vanishing first
$I$-vectors and the formula for $\Phi$, (C3) supplies the complete
original \texttt{CurveSystemVanishing} record.
% END CRT5

% BEGIN CRT6
\paragraph{CRT6. Recovering the gauge and its exact coefficient weight.}
Use $y_b$ for the opposite-side variables and define
\[
 F=X.hM(g,a,i;e,v')-p^{-E}\sum_{b,j}\Phi_{bj}z_j^{p^i}y_b .
\]
The gauge has zero restriction to $y=0$: every original monomial
in the first side alone has zero coefficient by $X.hL_{\rm axes}$,
and every other monomial has a positive-degree factor in the opposite
variables after the prescribed linear substitution. The displayed
main term vanishes on that axis as well, including its constant term.

Write $H'_{\kappa b}=\sum_k H_{\kappa k}v'_{bk}$. For each $b$ the
chain rule and the original differential condition $X.C_{\rm eq}$
give
\[
 \partial_{y_b}F=\Gamma_b(L,L')+
 \sum_{\kappa}H'_{\kappa b}
 \left(Z_\kappa-p^{-E}\sum_jT_{j\kappa}z_j^{p^i}\right).       \tag{C4}
\]
Here $L'$ are the coordinates of the opposite group pulled back along
$v'$, and $\Gamma_b$ is a power series with coefficients in $R$.
To check this on the actual leaf, the chain rule gives
$\sum_kv'_{bk}\partial_{(1-g,k)}hL$ followed by pullback.
For each $k$ with a nonzero multiplier the type is $a+i$, and
$(a+i-a).\mathrm{val}=i$ is odd. The differential condition expresses
this derivative as the sum of the leaf projections of the pairing
with $S.\mathrm{ze}(t)$ and an element of $\mathrm{IntGerm}(V,t)$.
All coordinates outside that leaf pull back to zero by the purity of
$e,v'$. Thus the projections may be removed after this substitution.
By definition the integral germ is obtained by substituting the
chosen two systems of coordinates into a series over $R$; summing
these series with the integral multipliers $v'_{bk}$ yields $\Gamma_b$.
Differentiating the displayed main term and using CRT4 gives (C4).
The type rule for $H$ and purity of $v'$ force $H'_{\kappa b}=0$
if $\kappa$ is outside the range of (C3).

Set $r_0=1/(p-1)$, $r_R=\rho+1/2-1/(8p^i)$. At right-variable
weight $r_0$, each opposite-side coordinate $L'_c$ has weight at
least $r_0$. Its recursion is (C1) with an integral linear term,
because $v'$ and the inverse opposite tangent basis are integral.
The assertion is preserved by that recursion:
$-r+p^r r_0\ge r_0$ for $r\ge1$, by $p^r-1\ge r(p-1)$.
Starting with zero and stabilizing degree by degree proves it.
With variable weights $(r_L,r_0)$, every coordinate in
$\Gamma_b$ has nonnegative weight and zero constant coefficient.
Each coefficient of the substitution is a finite sum of products,
so (C3)--(C4) imply $w_{r_L,r_0}(\partial_{y_b}F)\ge0$.

If $[z^\lambda y^\nu]F\ne0$, choose $b$ with $B=\nu_b>0$
by the axis result. The preceding derivative bound gives
\[
 v([z^\lambda y^\nu]F)+r_L|\lambda|+r_R|\nu|
 \ \ge\ r_0+(r_R-r_0)|\nu|-v_p(B).
\]
The restriction of the valuation to positive integers is the usual
$p$-adic exponent: integers are integral, and an integer prime to $p$
is a unit by B\'ezout and $v(p)=1$. If $v_p(B)=k$, then
$B\ge5^k\ge1+4k$, hence $k\le 4B/15$. On the other hand
\[
 r_R-r_0=\tfrac12-\frac{p}{p^2-1}-\frac1{8p^i}
 \ \ge\ \tfrac12-\tfrac5{24}-\tfrac1{40}=\tfrac4{15}.
\]
Since $B\le|\nu|$, the required original lower bound $r_0$ follows.
The weights here are exactly $r_1(p,i)-\delta(p,i)$ and
$r_2(p)-\delta(p,i)$. The coefficientwise bound, the exact first-axis
vanishing, and CRT4 give every field of \texttt{MConcl} on $U,U'$.
Together with CRT5 this proves the original \texttt{LemmaM}.
% END CRT6

% BEGIN CRT7
\paragraph{CRT7. The level-one refinement on these same rows.}
At $i=1$ let $z_c^*=\sum_{j:u_j=c}z_j$. Because the selected
indices are distinct, each $z_c^*$ is either zero or one left
variable. Modulo all monomials of total left degree at least $p^2$,
the first equation (C1) admits $L_c=z_c^*$ as a solution.
Indeed terms with $r\ge2$ disappear, and for $r=1$ the only
nonzero $z_c^*$ have $\alpha_{c1b}=0$ for every $b$ by goodness.
Uniqueness by the degree recursion shows these are the actual left
coordinates through that truncation. The second equation therefore gives
\[
 Z_\kappa\equiv-\sum_j\beta_{u_j,1,\kappa} z_j^p
                              \pmod{(z)^{p^2}}.
\]
Since $T_{j\kappa}=-p\beta_{u_j,1,\kappa}$ and $E=1$, the
$Z$-remainder of (C4) vanishes in this range. The opposite
coordinates are linear with integral coefficients modulo $(y)^p$,
since every nonlinear term in (C1) has total degree at least $p$.
Consequently $\Gamma_b(L,L')$ has integral coefficients throughout
$|\lambda|<p^2$, $|\nu|<p$: the two truncated coordinate families
are integral there, and multiplication and substitution preserve this
downward-closed range. If $|\nu|=0$ the axis coefficient vanishes.
Otherwise choose $0<\nu_b\le|\nu|<p$. Equation (C4) at
$(\lambda,\nu-\mathbf e_b)$ now shows that
$\nu_b [z^\lambda y^\nu]F$ is integral. Dividing by this unit gives
exactly \texttt{MPlusConcl}. All the witnesses from CRT4--CRT6
work simultaneously; this proves the original \texttt{LemmaMPlus}
for $1<d$.
% END CRT7

% BEGIN CRT8
\paragraph{CRT8. Attainment for the original odd block with either flag.}
For the window application take $n\ge1$ and prescribe the two
injective lists of $n$ letters in the indicated types. Their
indicator rows are saturated, the transpose indicators being integral
right inverses. CRT6 supplies a normal form on these very rows.
DJW15 converts it to the actual restricted perturbation and both
weighted row hypotheses. DJW20 then realizes the scalar comparison
by integral germs on the prescribed determinant letters. DJW21
gives, for the original product of the two letter determinants with
exponents $p^iy,y$, a system and coordinates fixed before the window
origin, with order zero admissible and with some admissible $y$ in
every interval $[Y_0,Y_0+S_0(p,i,n))$. When the selected rank is $n_0=0$, take
$y=Y_0$ as in DJW21. For $y>0$ the bound on the original valuation is
\[
 ny\left((p^i+1)\frac{p}{p^2-1}-1\right)
 +n(\log_p n+1)+\frac{(i+1)(j+j')(p^i+1)y}{p-1}.
\]
At $y=0$ the actual value is one. Here
$S_0=\lceil64n^2(i+3+3\log_p(n+1))\rceil$.
Both flags admit the required lists: the type $a$ on side $g$ and
the type $a+i$ on side $1-g$ are big, hence each has rank $n+j$.
Take the first $n$ entries of the respective supplied enumerations.
Under the variable renaming the leading minor is precisely the
selected-letter determinant, by its signed permutation expansion.
The same argument with the second flag interface yields
\texttt{BiSupply} and \texttt{RF.BiSupply} on their original data.

This section consumes the Cartier and row hypotheses of $X$; in
particular $R1$, $R2$, the integral splitting of the pairing, and
the integral-germ assertion in $C_{\rm eq}$ were used at their
stated types. It does not construct the global geometric family
giving rise to $X$.
% END CRT8
